% =============================================================================
% BAB III METODE -- teks panduan templat Sain-Tek-Kes 2026
% =============================================================================

\chapter{METODE}

Bab ini dapat diawali dengan kerangka pendekatan studi. Metode dapat berupa
percobaan laboratorium, percobaan lapangan, dan survei lapangan yang dirancang
sesuai dengan tujuan atau jenis penelitian, seperti: eksploratif, deskriptif,
koreksional, kausal, komparatif, eksperimen, tindakan
(\textit{action research}), pemodelan, analisis suatu teori, atau kombinasi
dari berbagai jenis penelitian. Judul Subbab pada Metode disesuaikan dengan
bidang keilmuan. Untuk tugas akhir yang menggunakan metode kualitatif, jelaskan
pendekatan yang digunakan, proses pengumpulan dan analisis informasi, dan
proses penafsiran hasil tugas akhir. Maksud dari perincian ini ialah untuk
menjamin keterulangan hasil. Berikut contoh subbab metode penelitian.

\section[Waktu dan Tempat]{Waktu dan Tempat (Judul subbab disesuaikan dengan bidang keilmuan)}

Deskripsi\dotfill

\noindent\dotfill

\noindent\dotfill

\section[Alat dan Bahan]{Alat dan Bahan (Judul subbab disesuaikan dengan bidang keilmuan)}

Deskripsi\dotfill

\noindent\dotfill

\noindent\dotfill

\section[Prosedur Kerja]{Prosedur Kerja (Judul subbab disesuaikan dengan bidang keilmuan)}

\subsection{Judul Sub-subbab (Kata dalam judul diawali huruf kapital dan dicetak tidak tebal)}

\begin{ipbsubsubcontent}
Berikut adalah contoh uraian dengan deskripsi pada sub-subbab. Pada sub-subbab
ini posisi paragraf lebih menjorok 0,5 cm dari paragraf di subbab. \dotfill

\noindent\dotfill

\noindent\dotfill

\noindent\dotfill

\noindent\dotfill

\noindent\dotfill

\noindent Contoh penulisan formula

\begin{equation*}
  f(x)=a_{0}+\sum_{n=1}^{\infty}
  \left(a_{n}\cos\frac{n\pi x}{L}
  +b_{n}\sin\frac{n\pi x}{L}\right)
\end{equation*}

\noindent Keterangan:

\noindent Berikan deskripsi pada setiap notasi formula di atas.
\end{ipbsubsubcontent}

\section{Analisis Data}

Deskripsi\dotfill

\noindent\dotfill

\noindent\dotfill
